\documentclass[11pt,a4paper]{article}

\usepackage[utf8]{inputenc}
\usepackage[T1]{fontenc}
\usepackage[margin=1in]{geometry}
\usepackage{booktabs}
\usepackage{amsmath}
\usepackage[hidelinks]{hyperref}
\usepackage{xcolor}
\usepackage{caption}

%% Numbers are GENERATED, never typed. Regenerate with:
%%   python3 paper/compute_numbers.py --latex
%% If numbers.tex is absent or has fallen behind, stop with one readable error
%% instead of letting every figure fail separately as an undefined control
%% sequence. Overleaf users hand-carrying files hit exactly this.
\IfFileExists{numbers.tex}{% GENERATED FILE - DO NOT EDIT.
% Produced by paper/compute_numbers.py from data/*.csv.
% Any manuscript number that is typed by hand is a bug: it will
% silently go stale as the ledger grows. That is precisely the
% failure mode this paper documents.

\newcommand{\NumbersStamp}{2026-09-19, 97 values}

\newcommand{\AbsMean}{8.8}
\newcommand{\AbsMedian}{2.0}
\newcommand{\AbsN}{47}
\newcommand{\AbsOthMeanRatio}{2.0}
\newcommand{\BSS}{+0.038}
\newcommand{\BSSHi}{+0.093}
\newcommand{\BSSLo}{-0.068}
\newcommand{\BSSWalk}{+0.046}
\newcommand{\BaseRate}{0.67}
\newcommand{\BeatAllHi}{0.390}
\newcommand{\BeatAllK}{186}
\newcommand{\BeatAllLo}{0.309}
\newcommand{\BeatAllN}{534}
\newcommand{\BeatAllP}{1.0}
\newcommand{\BeatAllRate}{0.348}
\newcommand{\BeatDriftK}{356}
\newcommand{\BeatDriftN}{534}
\newcommand{\BeatDriftP}{5.3\times10^{-15}}
\newcommand{\BeatDriftRate}{0.667}
\newcommand{\BeatMomK}{422}
\newcommand{\BeatMomN}{534}
\newcommand{\BeatMomP}{1.3\times10^{-43}}
\newcommand{\BeatMomRate}{0.790}
\newcommand{\BeatNaiveK}{345}
\newcommand{\BeatNaiveN}{565}
\newcommand{\BeatNaiveP}{8.1\times10^{-8}}
\newcommand{\BeatNaiveRate}{0.611}
\newcommand{\BrierBase}{0.2220}
\newcommand{\BrierModel}{0.2135}
\newcommand{\BrierN}{535}
\newcommand{\ConfirmN}{44}
\newcommand{\CorrAtwoAthree}{0.88}
\newcommand{\CovHi}{0.913}
\newcommand{\CovK}{503}
\newcommand{\CovLo}{0.862}
\newcommand{\CovN}{565}
\newcommand{\CovP}{6.7\times10^{-9}}
\newcommand{\CovRate}{0.890}
\newcommand{\DayAllK}{10}
\newcommand{\DayAllN}{38}
\newcommand{\DayAllP}{1.0}
\newcommand{\DayDriftK}{37}
\newcommand{\DayDriftN}{41}
\newcommand{\DayDriftP}{5.1\times10^{-8}}
\newcommand{\DayMomK}{39}
\newcommand{\DayMomN}{40}
\newcommand{\DayMomP}{3.7\times10^{-11}}
\newcommand{\DayNaiveK}{31}
\newcommand{\DayNaiveN}{41}
\newcommand{\DayNaiveP}{7.3\times10^{-4}}
\newcommand{\DirtyN}{6}
\newcommand{\EffFav}{11}
\newcommand{\EffFavShare}{0.478}
\newcommand{\EffHi}{0.670}
\newcommand{\EffLo}{0.292}
\newcommand{\EffN}{23}
\newcommand{\EffNeutral}{47}
\newcommand{\EffP}{1.00}
\newcommand{\EffUnfav}{12}
\newcommand{\ErrAgent}{95}
\newcommand{\ErrAgentShare}{0.785}
\newcommand{\ErrHuman}{26}
\newcommand{\ErrN}{121}
\newcommand{\FirstDay}{2026-08-01}
\newcommand{\KNineOneLat}{18}
\newcommand{\KNineThreeLat}{35}
\newcommand{\KOneOneEightLat}{1}
\newcommand{\LastDay}{2026-09-19}
\newcommand{\LatGeSeven}{35}
\newcommand{\LatGeThirty}{9}
\newcommand{\LatMax}{47}
\newcommand{\LatMaxClass}{data}
\newcommand{\LatMaxEvent}{2026-08-01}
\newcommand{\LatMaxRecord}{K117}
\newcommand{\LatMean}{6.1}
\newcommand{\LatMedian}{1}
\newcommand{\LatPninety}{19}
\newcommand{\LatPninetyfive}{32}
\newcommand{\LatPseventyfive}{8}
\newcommand{\MWP}{0.03}
\newcommand{\MWZ}{-2.14}
\newcommand{\NEff}{2.7}
\newcommand{\NumAssets}{7}
\newcommand{\NumDays}{49}
\newcommand{\NumScored}{565}
\newcommand{\NumWritten}{857}
\newcommand{\OthMean}{4.4}
\newcommand{\OthMedian}{1.0}
\newcommand{\OthN}{74}
\newcommand{\RepHi}{0.502}
\newcommand{\RepK}{50}
\newcommand{\RepLo}{0.330}
\newcommand{\RepRate}{0.413}
\newcommand{\RowsPerDay}{12.0}
\newcommand{\SkippedDays}{1}
\newcommand{\SpanDays}{50}
\newcommand{\VocabFreezeN}{77}
}{%
  \GenericError{}{numbers.tex not found}{}{%
    Run: python3 paper/compute_numbers.py --latex.
    On Overleaf, upload the whole paper/ directory, or link the project to the
    GitHub repository so generated files stay in step.}}
\providecommand{\NumbersStamp}{\textbf{MISSING}}
\captionsetup{font=small}

%% ---------------------------------------------------------------------------
%% AUTHOR DECISION — model disclosure.
%% The released dataset pseudonymises the forecasting model as model_A/model_B.
%% If you decide to name the model family in the paper, redefine this macro and
%% add a sentence to the Data Availability section. Leaving it as-is keeps the
%% manuscript consistent with the published CSVs.
\newcommand{\modelA}{\texttt{model\_A}}
\newcommand{\modelB}{\texttt{model\_B}}
%% ---------------------------------------------------------------------------

%% Title: kept to two lines. The subtitle's old tail ("Maintaining a Live
%% Measurement System") is now the first sentence of the abstract, where it
%% costs nothing. No manual \\ break -- let LaTeX choose it.
\title{Rules Are Not Mechanisms:\\
Operational Failure Telemetry from a Language-Model Agent}

\author{Abdullah Faruk Ciftler\\
\small ORCID: \href{https://orcid.org/0009-0001-9310-3812}{0009-0001-9310-3812}}

%% The generation stamp is printed on purpose: this is a living document, and
%% a numbers.tex that has fallen behind the data should be visible on page one
%% rather than only in CI.
\date{\today\\[2pt]\small Interim report, the ledger is ongoing\\[2pt]
\footnotesize Figures generated \NumbersStamp}

\begin{document}
\maketitle

\begin{abstract}
Agent evaluation is normally episodic: a task is attempted once and scored
pass or fail. This report covers a different setting. On \NumDays{}
forecast days between \FirstDay{} and \LastDay{} a language-model agent ran a live measurement system with real
consequences under human supervision: a personal forecasting ledger that writes
dated, sealed, pre-registered predictions and scores them against realised
outcomes. Alongside the forecasts we kept a structured log of every
operational error the agent made while running the process, recording the
error class, the mechanism that caught it, \emph{who} caught it, and how many
days it stayed hidden. We release both datasets.

The forecasting results are negative and were pre-registered. The agent beats
each of three baselines individually but beats all three simultaneously in only
$\BeatAllK/\BeatAllN = \BeatAllRate$ of scored rows ($p = \BeatAllP$,
one-sided), and on only \DayAllK{} of \DayAllN{} days when the forecast day
rather than the row is the unit. Prediction intervals are too wide
($\CovRate$ coverage against an $0.80$ target, exact binomial $p = \CovP$).
The up-move probabilities score about the same as a constant forecast at the
sample base rate: the Brier skill score is $\BSS$, and resampling whole days
puts its 95\% interval between $\BSSLo$ and $\BSSHi$.

The error log is the more unusual object. Across \ErrN{} records, detection latency
is right-skewed. The median is \LatMedian{} day, the mean \LatMean, the
maximum \LatMax{} days. Of the \ErrN{} records, \RepK{} ($\RepRate$) repeat
an earlier one. We name and
illustrate a recurring failure mode of agent-maintained systems: the agent
writes a correct rule into its own documentation and then does not follow it,
because \emph{documentation is not a mechanism}. Two clear instances and three
adjacent ones are described; the clearest went unenforced for \KNineThreeLat{}
days. The newest adjacent case runs the other way. A check was built for a
specific failure and caught it in the ledger, and hours later the same failure
reached the operator through conversation, a path the check does not read.
During the audit of this manuscript the pattern recurred once more, inside the
error log itself, and we report that instance too. Finally, we state the limitation that
governs all of this: the log is written by the same agent that made the errors,
so errors it never noticed are absent by construction. We propose a
capture--recapture design that would make the unobserved population estimable.
\end{abstract}

\section{Introduction}

Most published evaluation of language-model agents is \emph{episodic}. A
benchmark presents a task, the agent attempts it, and the attempt is scored.
This design answers ``can the agent do X?'' It does not answer a question that
matters more once agents are deployed: \emph{what happens to a system an agent
maintains over months, and how do its failures actually surface?}

Answering that requires a task benchmarks usually cannot provide. It has to
run continuously rather than once, so errors have time to compound or to be
caught, and it has to carry consequences, so the agent has something to be
wrong about. It also needs a human in the loop. Without someone who
occasionally notices what the agent missed, self-report cannot be checked at
all.

This paper reports on such a task and releases the resulting data. The agent
maintains a forecasting ledger over \NumAssets{} financial assets, six of them
forecastable. Each business day it
records point forecasts, $80\%$ intervals and up-move probabilities at 1-day,
1-week and 1-month horizons, before outcomes are known. The ledger seals them
and scores them automatically. The analysis plan was pre-registered, and
protocol changes increment a version so that observations produced under
different rules never land in the same pool.

Our contribution is not a forecasting method. On the contrary, the forecasting
results are null, and we report them as such. The contribution is the
\emph{operational telemetry} collected alongside: a structured record of every
mistake the agent made while running the process, with detection latency and
attribution.

\paragraph{Contributions.}
\begin{enumerate}
  \item A released, anonymised, longitudinal dataset of language-model agent
        operational errors with detection mechanism, detecting party, and
        detection latency in days (\S\ref{sec:errors}).
  \item Pre-registered null results on short-horizon forecasting, including
        interval mis-calibration in a specific and correctable direction
        (\S\ref{sec:forecast}).
  \item A named failure mode, \emph{rules written but not mechanised},
        illustrated by two clear instances and three adjacent ones, with a
        measured repeat rate (\S\ref{sec:mechanism}).
  \item A design proposal that would make the unobserved error population
        estimable rather than merely acknowledged (\S\ref{sec:limits}).
\end{enumerate}

\section{Related work}\label{sec:related}

\paragraph{Agent evaluation is episodic.} Benchmarks such as
SWE-bench~\cite{jimenez2024swebench} present a task, let an agent attempt it,
and score the attempt. The design is well suited to measuring capability and
poorly suited to measuring what happens to a system an agent tends over time.
An agent that scores well on isolated tasks may still degrade a system it
maintains, and the degradation is invisible to a benchmark that never runs
twice on the same artefact.

Continuing deployments do exist and are reported. Williams et
al.~\cite{williams2026pomona} ran Pomona, an agent that opens small repair
pull requests, for three months inside Bloomberg and evaluated it with a
survey of senior engineers. The framing is close to ours in one respect and
opposite in another. Close, because their stated aims are reducing technical
debt and keeping engineer trust, which are the two things a maintained system
runs on. Opposite, because they measure what the agent produced and we measure
what it got wrong. Their unit is the accepted pull request; ours is the logged
mistake, its detection latency, and who caught it.

\paragraph{The human in the loop is a cost, not a free correction.} Our
process has a supervisor who catches what the agent does not, and
\ErrHuman{} of \ErrN{} records were found that way. It is tempting to read
that as reassurance. Garousi~\cite{garousi2026oversight} argues the opposite
and we think he is right: reviewing and reworking what an agent produces is
not optional and not free, and the volume of suggestions imposes a cognitive
load of its own. A supervision rate is therefore a cost figure as much as a
safety figure, and an agent that needs less of it is doing better than one
whose errors merely get caught. We cannot yet separate the two, for the reason
given in \S\ref{sec:limits}: we never recorded whether the agent would have
caught what the human caught.

\paragraph{Forecast scoring.} The scoring apparatus used here is standard.
Proper scoring rules and their decomposition are treated by Gneiting and
Raftery~\cite{gneiting2007scoring}; the probability score itself dates to
Brier~\cite{brier1950}. Our contribution to this literature is nil: we use it
as instrumentation and report that it returns a null. The practice of scoring
dated judgements against baselines, and of taking the comparison rather than
the raw hit rate as the unit of evidence, follows the forecasting-tournament
tradition described by Tetlock and Gardner~\cite{tetlock2015}.

\paragraph{Pre-registration.} Fixing the analysis plan before the data is seen
is the standard defence against selective reporting~\cite{nosek2018prereg}. It
is normally applied to a study; here it is applied to a running process, which
raises a question the literature does not settle: what counts as a
pre-registration when data accrues daily and hypotheses are added as the
process reveals them. We version the protocol, split the pool, and mark
anything added after the fact as exploratory. That is a convention we adopted,
and we make no claim that it is the right one.

\paragraph{Latent conditions and incident analysis.} The pattern we name in
\S\ref{sec:mechanism} is a close relative of Reason's \emph{latent
conditions}~\cite{reason1990}: a defect introduced long before the failure it
enables, dormant because nothing exercises it. The difference is the author.
In Reason's account latent conditions are left by designers for operators to
trip over; here the agent is both, and writes the very rule it will later fail
to apply. Whether this makes the failure more or less tractable than the
human case is open; we can only report that it happened repeatedly and, once,
inside the audit of this manuscript (\S\ref{sec:limits}).

\paragraph{Technical debt.} The mechanism of \S\ref{sec:mechanism} has an
older name in software engineering. Behutiye et al.~\cite{behutiye2017td},
reviewing technical debt in agile development, report that among the
management strategies found in the literature the two that recur most are
refactoring and \emph{increasing the visibility of the debt}, and that
pressure for quick delivery is among its most common causes. Both map onto
what we see. A rule that is written down but never wired into the process is
debt of exactly this kind: the cost is deferred, it accrues quietly, and it is
paid in repeat failures. Our latency field is a visibility instrument in their
sense, since it measures how long a defect survived unnoticed rather than how
many defects there were, and the causal note is ours too, as several of our
records begin with a warning that was on screen and read too fast. What this
setting adds is that debtor and creditor are the same agent, which takes on
the debt in the very act of writing the lesson down.

The Dagstuhl manifesto on reframing technical
debt~\cite{avgeriou2025reframing} names tooling and data collection among the
places current practice falls short, and our records are a small instance of
both shortfalls at once. The rule that went unenforced was written in a
document no tool reads, and the only reason its cost is visible here is that
we happened to be keeping a log with a latency field in it. Nothing about the
process would have surfaced it otherwise, and nothing in the tooling we were
using was looking.

\paragraph{What we did not find.} Longitudinal agent deployments are
reported, as \cite{williams2026pomona} shows. What we did not find is a
released dataset from one of them recording the agent's own operational
failures with detection latency and attribution, nor work applying
capture--recapture to estimate an agent's unobserved error population. The
distinction matters and is narrower than the claim we first wrote: the gap is
not that nobody runs agents for months, it is that what those months produce
in the way of failure data stays inside the organisation that collected it.
If such a release exists we would like to know, and we say this to invite
correction rather than to claim novelty by absence of search.

\section{Setup}\label{sec:setup}

\paragraph{The task.} A daily process over \NumAssets{} assets (six
forecastable, pseudonymised \texttt{A1}--\texttt{A7}): collect prices from machine-readable
sources, record a market-state snapshot, resolve any forecasts that have come
due, scan the day's news, write new forecasts, and commit. A fixed engine performs the steps that must be
deterministic, and the agent performs the ones that need judgement. The engine
never fetches data and never writes a forecast. The agent never computes a
score. The split is deliberate. Arithmetic must not hallucinate, and neither
must data collection.

\paragraph{What is recorded.} Every forecast carries a point estimate, an
$80\%$ interval, an up-move probability, a stated confidence, the evidence
items it was derived from, and an invalidating condition. Resolution is
automatic against the realised value and against three baselines: \emph{naive}
(no change), \emph{drift} (mean of the last five observed changes) and
\emph{momentum} (the last observed change). All three use only information
available at forecast time.

\paragraph{Pre-registration.} The analysis plan fixes, in advance, which
hypothesis is primary, which test is applied, the minimum $n$, and the
reporting order. Findings that are not in the plan are labelled exploratory
and are not counted as evidence. Rule changes increment a protocol version and
split the pool.

\paragraph{Scale.} \NumDays{} forecast days, \NumWritten{} forecasts
written, \NumScored{} scored, \ErrN{} error records. The window runs from
\FirstDay{} to \LastDay, which is \SpanDays{} calendar days, and on
\SkippedDays{} of them the process was not run at all. Prices for a missed day
can be fetched afterwards. The news scan cannot, so whatever moved the market
that day is absent from the record for good. This is an interim report; the
ledger continues.

\paragraph{Excluded rows.} The resolver looks forward up to six days when a
target date has no price, which is correct for holidays and for fund prices
published a day late. It cannot tell a closed market from a day nobody
recorded, and on the missed day it would have scored two days of movement as
one. Of the resolved rows, \DirtyN{} took their price this way on an open
market day. They remain in the released data, flagged, and are excluded from
every figure below.

\section{Forecast results are null}\label{sec:forecast}

\begin{table}[t]
\centering
\small
\begin{tabular}{lrrrrr}
\toprule
Baseline & rows beaten & rate & $p$ (rows) & days won & $p$ (days) \\
\midrule
naive    & \BeatNaiveK/\BeatNaiveN & \BeatNaiveRate & $\BeatNaiveP$
         & \DayNaiveK/\DayNaiveN & $\DayNaiveP$ \\
drift    & \BeatDriftK/\BeatDriftN & \BeatDriftRate & $\BeatDriftP$
         & \DayDriftK/\DayDriftN & $\DayDriftP$ \\
momentum & \BeatMomK/\BeatMomN & \BeatMomRate & $\BeatMomP$
         & \DayMomK/\DayMomN & $\DayMomP$ \\
\textbf{all three at once} & \textbf{\BeatAllK/\BeatAllN} &
\textbf{\BeatAllRate} & \textbf{\BeatAllP} & \textbf{\DayAllK/\DayAllN} &
\textbf{\DayAllP} \\
\bottomrule
\end{tabular}
\caption{Beating any single baseline is easy in a trending sample. Beating all
three at once is the number that matters, and it fails. The day columns treat
each forecast day as one observation: a day is won when more than half of its
rows beat the baseline, and a day split exactly in half drops out. All
$p$-values are one-sided sign tests. Wilson 95\% CI for the joint row:
$(\BeatAllLo,\,\BeatAllHi)$.}
\label{tab:baselines}
\end{table}

Table~\ref{tab:baselines} states the central negative result. Report only the
naive comparison and the headline writes itself: the agent beats the baseline,
$p = \BeatNaiveP$. It would have been an artefact. A fraction $\BaseRate$ of
realised moves in the window were upward, so any small positive point forecast
beats a zero-change baseline mechanically, and a headline built on that
measures the market rather than the forecaster. The joint criterion was fixed
in advance for exactly this reason.

The row-level $p$-values treat every row as independent, and they are not:
one day's forecasts share the same news and move with the same markets. The
day columns of Table~\ref{tab:baselines} use the day as the unit instead,
which the pre-registration requires. The conclusion does not move. Every single
baseline is still beaten, and the joint criterion is met on \DayAllK{} of
\DayAllN{} days.

\paragraph{Intervals are mis-calibrated in a correctable direction.}
Coverage is $\CovK/\CovN = \CovRate$ (Wilson 95\% CI
$\CovLo$--$\CovHi$) against an $0.80$ target; the exact binomial probability
of coverage this high or higher under a true $0.80$ is $\CovP$. The intervals are too wide. Unlike
the point forecasts, this is a measurable and fixable defect, and it is where
the ledger's own ongoing work is directed.

\paragraph{Probabilities show no demonstrated skill.} Mean Brier score is
$\BrierModel$ ($n=\BrierN$) against $\BrierBase$ for a constant forecast at the
sample base rate of $\BaseRate$, a Brier skill score of $\BSS$. When this
report was first drafted the score was negative. It is positive now, and that
is not evidence of anything. Resampling whole forecast days gives a 95\%
interval from $\BSSLo$ to $\BSSHi$, which contains zero with room to spare. A
score that changes sign between two versions of the same report is measuring
the sample.

There are two defensible base rates, and we lead with the one the
pre-registration did not name. The plan specifies a walk-forward rate that
uses only rows already resolved; \texttt{reproduce.py} computes it and returns
$\BSSWalk$. The figure above uses the full-sample rate instead, a look-ahead no
real forecaster would have had, because it gives the lower of the two numbers.
Picking the friendlier one after seeing both would be a correction in our own
favour, the bias \S\ref{sec:errors} tests for. The gap between them is much
smaller than the width of the interval, so the choice changes no conclusion.

\paragraph{Power.} These figures overstate the evidence. The assets move in
correlated blocks: the two foreign-equity funds have a daily return
correlation of $\CorrAtwoAthree$, and two precious-metal instruments track the
same underlying. A Kish-style effective count,
$k^{2}/\sum_{ij}\rho_{ij}$, gives $\NEff$ effectively independent series per
day rather than the nominal \NumAssets. Clustering by day, as in
Table~\ref{tab:baselines}, absorbs this. It does not absorb overlap between
days: one-week forecasts written on consecutive days share most of their
target window, so neighbouring days are not independent either.
With \NumDays{} days, no forecasting conclusion here should be read as more
than ``not demonstrated''.

\section{Operational failure telemetry}\label{sec:errors}

Every operational mistake is logged with a controlled vocabulary: error class
(8 values), the mechanism that caught it (10 values), the detecting party
(agent or human), the direction of the resulting correction, and the number of
days between the event and its detection.

\begin{table}[t]
\centering
\small
\begin{tabular}{lr}
\toprule
records & \ErrN \\
detected by agent / human & \ErrAgent{} / \ErrHuman{} \quad (agent share $\ErrAgentShare$) \\
detection latency: median / mean / max & \LatMedian{} / \LatMean{} / \LatMax{} days \\
latency $p_{75}$ / $p_{90}$ / $p_{95}$ & \LatPseventyfive{} / \LatPninety{} / \LatPninetyfive{} days \\
latency $\geq 7$ days & \LatGeSeven/\ErrN \\
latency $\geq 30$ days & \LatGeThirty/\ErrN \\
repeat of an earlier record & \RepK/\ErrN{} $= \RepRate$ \; CI95 $(\RepLo,\,\RepHi)$ \\
\bottomrule
\end{tabular}
\caption{Error log summary. Latency is heavily right-skewed: most errors are
caught within a day, a minority survive it for weeks.}
\label{tab:errors}
\end{table}

\paragraph{Two features stand out.} First, latency is skewed: the median is
\LatMedian{} day while the mean is \LatMean{} and the maximum is \LatMax{}
days. Routine cross-checking catches most errors within a day; a minority
survive it for weeks. The tail is bounded by the observation window, and
the longest-hidden record shows how tightly: \LatMaxRecord{} has an event
date of \LatMaxEvent, and the ledger's first forecast day is \FirstDay. The
maximum says
how long one error did hide, and nothing about how long an error \emph{can}
hide. Second, \RepK{} of \ErrN{} records ($\RepRate$) repeat an earlier record.
The agent returns to failure classes it has already documented.

\paragraph{The longest-hidden error was usually right.} \LatMaxRecord{} is a
\LatMaxClass{} error, not an omission. The daily run stored the domestic
index level at the start of the session, so the stored value was an intraday
reading, and the day-to-day change built from it ran morning to morning rather
than close to close. On most days the two agree, because the index has barely
moved by the time the run starts. They part on days with a large move after
the reading, and those are the days the field exists for. Each correct day
made the next one less likely to be checked. It surfaced when a jump alarm
fired on a real move and verifying the move exposed the reading time. Nothing
was rewritten backwards: those were the values actually read on those
mornings, and a separately named previous-close field now sits beside them.

\paragraph{Self-correction appears directionally unbiased.} Of the records
whose effect on the ledger's own scores was classified, \EffFav{} were
favourable to the agent and \EffUnfav{} unfavourable (two-sided exact
$p = \EffP$; favourable share $\EffFavShare$, CI95 $\EffLo$--$\EffHi$). This is the one result here where a
\emph{significant} finding would have been bad news: a systematic skew toward
self-favouring corrections would indicate that corrections are chosen after
seeing which way they cut. No such skew is detected, though the test is
underpowered.

\section{Rules are not mechanisms}\label{sec:mechanism}

One recurring pattern is specific enough to name. The agent diagnosed a real
problem correctly, wrote the correct rule into its own documentation, and then
did not apply it. Repeatedly. A document is not an enforcement point.

We are deliberate about the count. \textbf{Two} records are clear instances: a
documented rule existed, no code enforced it, and the omission had a cost.
\textbf{Three} further records are adjacent but weaker. One is a missing
validation rather than an unfollowed rule, one is a wrong causal diagnosis
that was later corrected, and one is a mechanism that existed and did not
cover the path the failure took. We list all five and mark which is which,
rather than reporting ``five instances'' of a category stretched to fit.

\begin{table}[h]
\centering
\small
\begin{tabular}{llr p{5.6cm}}
\toprule
Record & Fit & Latency & The rule that existed but was not enforced \\
\midrule
\textbf{K93} & clear & \KNineThreeLat{} d & ``Index values must be cross-checked
against the previous close.'' Written after a $5\%$ index mis-read; no code
ever checked. \\
\textbf{K91} & clear & \KNineOneLat{} d & Three rules (shadow intervals, an
adversarial pre-commit review, a second ledger's snapshot) were documented.
None was in the daily runner. \\
K80 & adjacent & 0 d & Market-state fields silently disappeared between days.
A missing validation, not an unfollowed rule. \\
K87 & adjacent & 1 d & A prior diagnosis of a provider anomaly was itself wrong.
The behavioural rule survived; its stated cause did not. Included because it
shows a documented \emph{explanation} can be wrong while the \emph{rule} it
justifies stays useful. \\
K118 & adjacent & \KOneOneEightLat{} d & A check for index readings taken
before the session settled was committed to the ledger. Hours later the agent
quoted such a reading to the operator in conversation. The mechanism existed;
its scope was the ledger. \\
\bottomrule
\end{tabular}
\caption{Instances of the pattern, with an explicit judgement of fit. Two are
clear, three are adjacent and reported as such.}
\label{tab:mechanism}
\end{table}

The cost is measurable. A documented interval-recording rule was followed on
the day it was written, ignored for the next three days, then resumed when
someone happened to remember it. Those three days are gone. The quantity
cannot be reconstructed after the fact, so the hypothesis pool they were meant
to fill is permanently short by three observations.

\paragraph{A mechanism has a scope.} K118 belongs in the table because it runs
the other way. After an intraday index reading had been mistaken for a close,
the agent built the check the clear cases lacked, and the check worked on the
ledger. The same afternoon the agent told the operator an index had closed up
by an amount it had read before the final hour, and the actual close was
several times larger. The ledger never held the wrong number, so no ledger
check could have seen it. The record was logged against the operator's
interest, which is the direction a reader should worry about. Mechanising a
rule is necessary and still leaves a gap: a mechanism covers the path it is
wired into, and an agent that talks to people has paths no check reads.

\paragraph{Why this should generalise.} The pattern does not depend on the
particular agent or task. It depends on a structural asymmetry. An agent that
writes its own operating documentation is also the only reader of it, and
nothing enforces reading. A rule that survives only as prose therefore depends
on the agent remembering it at the moment it applies. That is precisely the
moment the agent is busy with whatever caused the rule to be written.

\paragraph{A failed exploratory test, and why it is weaker than it looks.}
We expected errors of \emph{absence} (a missing field, an unenforced rule) to
hide longer than errors of \emph{wrong value}, since a wrong number can be
cross-checked against another number and a missing one cannot. Absence-like
records ($n = \AbsN$) have median latency $\AbsMedian$ days and mean
$\AbsMean$; the remainder ($n = \OthN$) have median $\OthMedian$ and mean
$\OthMean$. A Mann--Whitney test gives $z = \MWZ$, two-sided $p = \MWP$. When this
manuscript was first drafted the same test returned nothing; the separation
appeared as records accumulated, and it now sits on the conventional side of
$0.05$ in the predicted direction, close enough to the line that a few more
records could move it back.

We report it, and we report why it is not evidence even though it
succeeded. The hypothesis was formed after seeing the latency distribution.
Worse, the operationalisation was too: ``absence-like'' is not a field in the
log but a grouping of two of the eight error classes, and we chose which two
\emph{after} looking. With eight classes there are $2^{8}$ possible groupings,
and nothing constrained our choice except the shape of the data. A test whose
hypothesis, grouping and direction are all selected post hoc has no
interpretable $p$-value, and this is exactly the situation in which a
$p$ below $0.05$ is least informative: the grouping was chosen to make the
separation visible. The number above describes this sample. It is not a test,
and we would have reported it identically had it come out null. The hypothesis
is now registered prospectively, with the grouping fixed in advance, and will
be answered by records this manuscript has not seen. The medians are
$\AbsMedian$ and $\OthMedian$ days and the means differ by a factor of
$\AbsOthMeanRatio$, so whatever separation exists is in the tail. The
longest-hidden record of all does not fit the hypothesis: \LatMaxRecord{} is a
wrong value, and it hid because it was usually right.

\section{Limitations}\label{sec:limits}

\paragraph{The log is self-reported.} It is written by the same agent that
made the errors. Errors the agent never noticed are absent by construction,
and no amount of internal analysis recovers them. The \ErrHuman{}
human-detected records hint at a miss rate without estimating one. We never
recorded whether the agent \emph{would} have caught what the human caught, so
the two detectors cannot be compared.

\paragraph{A design that would fix this.} Capture--recapture, standard in
epidemiology and software defect estimation, requires only that overlap be
recorded: for each error, whether each detector found it independently. With
overlap, the unobserved population becomes estimable rather than merely
acknowledged. We have added the field prospectively; it carries no data yet.

\paragraph{The pattern recurred during this audit, in the log itself.}
While auditing this manuscript we recomputed every latency value from the two
dates that define it. Three of \ErrN{} disagreed with the stored figure. Two
were off by a day. One, the record used above as the clearest example of the
pattern, stored \emph{four hundred} days where the arithmetic gives
\KNineThreeLat. The value had been entered by hand; the field is derivable; no check
existed. A four-hundred-day latency was impossible on its face, since the
ledger is \NumDays{} days old. Neither the agent that wrote the figure nor the
human who reviewed the release noticed.

The rule was already written. The project's own operating document states that
latency equals the logged date minus the event date. It was written down, it
was correct, and it was not enforced, so a derivable field was typed by hand
and never checked. The manuscript's most striking figure was itself a product
of the failure mode the manuscript describes.

The error reached the published dataset, this paper's abstract, and three
archived releases before it was caught. Those archives cannot be edited; the
correction is stated here and in the revision notes rather than made silently.
The check now exists and fails the daily run if any latency disagrees with its
dates. We report the episode because a paper about unenforced rules that
quietly fixed its own would be worth less than one that does not.

\paragraph{Self-categorisation, not just self-report.} The limitation above is
usually stated as ``the log is self-reported''. It is worse than that. The
agent also chose the error classes, designed the controlled vocabulary those
classes belong to, and decided which records count as instances of the pattern
in \S\ref{sec:mechanism}. Section~\ref{sec:mechanism} marks three of five
instances as merely adjacent for this reason, but the marking is itself the
agent's judgement. An independent re-classification of the released log would
be a useful contribution and we have not done one.

\paragraph{Single agent, single task, small $n$.} One agent, one task, one supervisor, \NumDays{} days. Nothing here establishes a rate that transfers to other
agents. The dataset's value is that it exists and is open, not that it is
representative.

\paragraph{Classification is partly retrospective.} The controlled vocabulary
was introduced after \VocabFreezeN{} records had accumulated; those were mapped to
it \emph{after} the data was visible. The released log therefore distinguishes
an exploratory pool from a confirmatory pool beginning at the record where the
vocabulary was fixed. Tests intended as evidence use the confirmatory pool
only, which currently holds \ConfirmN{} records.

\section{Data availability}

All data, documentation and analysis code are released under CC~BY~4.0 (data
and documentation) and MIT (code):

\begin{center}
\href{https://doi.org/10.5281/zenodo.22643370}{\texttt{doi:10.5281/zenodo.22643370}}
\quad
\href{https://github.com/farukciftler/forecast-discipline-ledger}{\texttt{github.com/farukciftler/forecast-discipline-ledger}}
\end{center}

The repository contains the pre-registration, methods, codebook, case
descriptions, and a dependency-free script that regenerates every number in
this paper from the released CSVs. Every figure and table above is produced by
\texttt{paper/compute\_numbers.py}.

The exporter that produces the anonymised release and the independent auditor
that verifies it are deliberately \emph{not} published: the exporter carries
the alias mapping and the auditor carries the forbidden-term list, so
publishing either would defeat the anonymisation.

\paragraph{What is withheld.} Amounts, balances, unit counts, price levels,
institution, instrument, vendor and person names, and free text. Assets appear
as \texttt{A1}--\texttt{A7}, the forecasting models as \modelA{} and \modelB{},
and country-specific market variables under role labels; global public
benchmarks keep their usual names. The accepted residual risk is that a fund's
return series could in principle be matched against public NAV series and the
instrument identified, which reveals nothing about any person.

\paragraph{Use of AI in preparing this manuscript.} The system under study
is a language-model agent, and this manuscript was itself drafted with
language-model assistance under the author's direction. The distinction that
matters is between the prose and the numbers. Every figure in the text is
generated from the released data by \texttt{compute\_numbers.py} and checked
by \texttt{check\_paper.py}, which fails the build if a number is typed by
hand rather than derived; the audit in \S\ref{sec:limits} describes what
happened the one time that guard did not yet exist. The author read the data
behind every claim and is responsible for the content, including the errors.
We state this because a paper about an agent's operational failures would be a
poor place to be quiet about it.

\section{Conclusion}

An agent ran a real measurement system on \NumDays{} days and produced
forecasts with no demonstrated skill. That half of the result is dull. It was
also pre-registered, which is the only reason it could not later be reframed
as something better.

What the process recorded about itself is more useful. Most errors surfaced
within a day, but \LatGeSeven{} of \ErrN{} took a week or longer and the
slowest took \LatMax{} days. A fraction $\RepRate$ of them repeat an earlier record. The
corrections the agent made to its own scores show no detectable bias in its own
favour, though that test is underpowered. The longest-hidden failure was a
field that was right on quiet days and wrong on the days that mattered. Several
of the other long ones were rules the agent had written down and then did not
apply, and when a mechanism was finally built for one failure, it guarded the
ledger and missed the conversation beside it.

Capability benchmarks answer how often an agent is right. An agent that
maintains something also needs its being wrong to become visible, and the
evidence here is that writing the lesson down does not make it so. Whether that
holds beyond one agent and \NumDays{} days is exactly what we cannot say from
this sample, which is why the data is released rather than summarised.

\bibliographystyle{plain}
\bibliography{refs}

\end{document}
